% C015 椭圆焦点三角形四心坐标性质——结论
% 本文件是正文片段，沿用宿主模板中的 statementbox 环境。
\begin{statementbox}
\textbf{C015\quad 椭圆焦点三角形四心坐标性质}

设椭圆
\[
 C:\ \frac{x^2}{a^2}+\frac{y^2}{b^2}=1,
 \qquad a>b>0,\quad c^2=a^2-b^2,\quad e=\frac ca\in(0,1),
\]
其焦点为 $F_1(-c,0)$、$F_2(c,0)$。取椭圆上的点
$P(x_0,y_0)$，且 $y_0\ne0$（即排除长轴的两个端点），
则 $\triangle PF_1F_2$ 是非退化三角形，其四心分别为：
\begin{enumerate}
 \item \textbf{重心（中线交点）}
 \[
 \boxed{G\left(\frac{x_0}{3},\frac{y_0}{3}\right)}.
 \]
 \item \textbf{内心（内角平分线交点）}
 \[
 \boxed{I\left(ex_0,\frac{cy_0}{a+c}\right)}.
 \]
 这里 $y_I$ 与 $y_0$ 同号，内切圆半径为
 $\displaystyle r=\frac{c|y_0|}{a+c}$。
 \item \textbf{垂心（三条高线交点）}
 \[
 \boxed{H\left(x_0,\frac{c^2-x_0^2}{y_0}\right)}.
 \]
 \item \textbf{外心（三边中垂线交点）}
 \[
 \boxed{O'\left(0,\frac{x_0^2+y_0^2-c^2}{2y_0}\right)}.
 \]
\end{enumerate}

\textbf{一眼记住四个突破口：}
重心看\textbf{平均}；内心看\textbf{切线长与面积}；
垂心看\textbf{竖直高线}；外心看\textbf{底边中垂线}。

\textbf{最短联动公式：}
\[
 \boxed{x_I=\frac{PF_1-PF_2}{2}=ex_0},\qquad
 \boxed{H=(x_0,y_0)-2O'}.
\]
后一式中的坐标向量都以椭圆中心 $O=(0,0)$ 为原点；
等价于三角形欧拉线公式 $\overrightarrow{O'G}=\tfrac13\overrightarrow{O'H}$。

\textbf{两条最有用的轨迹：}
若把 $G=(X,Y)$、$I=(X,Y)$ 分别视为动点，便有
\[
 \boxed{\frac{X^2}{a^2/9}+\frac{Y^2}{b^2/9}=1}
 \quad(\text{重心}),
\]
\[
 \boxed{\frac{X^2}{c^2}+
 \frac{Y^2}{\bigl(bc/(a+c)\bigr)^2}=1}
 \quad(\text{内心}).
\]
两条轨迹均\textbf{不包含横轴上的两个端点}，因为原三角形在 $y_0=0$ 时退化。
\end{statementbox}
